% ---------- Parte 1: diseño experimental y cómo se corrió ----------
\section{Diseño experimental}
\portadilla{1}{Diseño experimental: qué se midió y cómo se corrió}{Cuatro formas de correr las mismas pruebas, en dos máquinas, y por qué hizo falta cada una.}{Primera parte: cómo se hicieron las pruebas. Es la base de todo lo demás, porque un resultado vale lo que vale el experimento que lo produce.

Vamos a ver las máquinas, de dónde salieron los programas, cómo se corrieron con Docker y sin Docker, qué es el programa del equipo y qué se hizo en el i5-10400.}

\begin{frame}{Cuatro máquinas en juego: dos del paper y dos nuestras}
  \fuente{de Wolff et al., sección 6; datos/ultra9/equipo.txt; datos/i5\_10400/00\_entorno.txt}
  \footnotesize
  \begin{tabularx}{\linewidth}{@{}H{22mm}LLLL@{}}
    \toprule
    & \textbf{Paper, figura 5} & \textbf{Paper, figura 6} & \textbf{Réplica A} & \textbf{Réplica B}\\
    \midrule
    Procesador & 4 × Xeon E7-8867 & Xeon w7-3465X & Core Ultra 9 285 & Core i5-10400\\
    Núcleos e hilos & 72 núcleos, 144 hilos & 28 núcleos, 56 hilos & 24 núcleos (8 P + 16 E), sin SMT & 6 núcleos, 12 hilos (SMT)\\
    Chips y memoria & 4 sockets, 4 nodos NUMA & 1 socket & 1 socket y 1 nodo NUMA; L3 de 36 MB; RAM en 1 canal & 1 socket\\
    Sistema & — & — & Ubuntu 24.04, y Docker con Ubuntu 22.04 & WSL2 sobre Windows, Ubuntu 24.04\\
    Qué se midió & Rendimiento & Latencia & Todo: 4 tandas, controles, RAM, corrección & Corrección y una pasada rápida de \cod{bench.cpp}\\
    \bottomrule
  \end{tabularx}
  \vspace{2mm}
  \begin{idea}
    Ninguna de nuestras máquinas tiene 4 sockets. Por eso se comparan \textbf{órdenes y proporciones} entre algoritmos, no Mop/s absolutos entre máquinas.
  \end{idea}
  \note{Hay cuatro máquinas en juego. El paper usó dos: cuatro Xeon con 144 hilos para el rendimiento y un Xeon w7 de un solo chip para la latencia.

  Nosotros usamos otras dos. El Core Ultra 9 285 es un procesador híbrido de 24 núcleos: 8 rápidos, llamados P, y 16 eficientes, llamados E. Ahí se hizo la réplica completa. El i5-10400 tiene 6 núcleos con dos hilos cada uno y corre Linux dentro de Windows, con WSL2. Ahí se hicieron la corrección y una pasada rápida.

  La consecuencia importante: ninguna de las dos tiene cuatro chips. Así que no comparamos Mop/s entre máquinas, sino el orden de los algoritmos y cuánto rinde cada uno respecto a otro.}
\end{frame}

\begin{frame}{Qué se corrió en cada máquina y de qué forma}
  \fuente{EXPLICACION.md, Partes 0, 4 y 6.6; resultados/david-ultra9-285/LEEME.md}
  \scriptsize
  \setlength{\tabcolsep}{3pt}
  \begin{tabularx}{\linewidth}{@{}>{\raggedright\arraybackslash}p{36mm}CCCCCC@{}}
    \toprule
    & \multicolumn{5}{c}{\textbf{Core Ultra 9 285}} & \textbf{i5-10400}\\
    \cmidrule(lr){2-6}\cmidrule(l){7-7}
    & Docker, binarios oficiales & Sin Docker, binarios oficiales & Sin Docker, \cod{bench.cpp} & Sin Docker, \cod{bench.cpp} alineado & Pasada rápida de \cod{bench.cpp} & Pasada rápida de \cod{bench.cpp}\\
    \midrule
    Rendimiento (figura 5) & ✓ & ✓ & ✓ & ✓ & ✓ & ✓\\
    Latencia (figura 6) & ✓ 5 rondas & ✓ 15 rondas & ✓ & — & ✓ & ✓\\
    Hilos fijados a núcleos & ✓ & ✓ & ✓ & ✓ & — & —\\
    Dos hilos en pares de núcleos & ✓ 3 rondas & ✓ 10 rondas & ✓ & — & — & —\\
    Memoria retenida & ✓ & ✓ & — & — & — & —\\
    Receta exacta, 10 s × 7 & ✓ 2 escenarios & ✓ 6 escenarios & — & — & — & —\\
    Tráfico real de RAM & ✓ & — & ✓ & ✓ & — & —\\
    ASan, TSan, boceto & ✓ rivales & — & ✓ rivales & — & ✓ & ✓\\
    \midrule
    Mediciones & 3\,720 & 4\,660 & 2\,976 & 2\,160 & 168 & 168\\
    \bottomrule
  \end{tabularx}
  \vspace{1mm}
  \begin{idea}La pasada rápida es el mismo guion, \cod{ejecutar\_pruebas.sh}, en las dos máquinas: es la comparación directa entre ellas.\end{idea}
  \note{Esta tabla es el mapa de toda la réplica. Las columnas son las formas de correr las pruebas; las filas, los experimentos.

  En el Ultra 9 hay cuatro tandas grandes: los binarios oficiales con Docker y sin Docker, el programa del equipo sin Docker, y una copia de ese programa con una corrección de una palabra que explicaremos en la parte 6.

  La penúltima columna es importante. El guion del equipo, ejecutar\_pruebas.sh, se corrió igual en las dos máquinas. Eso nos da una comparación directa entre el i5 y el Ultra 9 con el mismo programa y los mismos parámetros. La veremos en la parte 5.}
\end{frame}

\begin{frame}{Más de 14\,000 mediciones, ninguna fallida}
  \fuente{EXPLICACION.md, Parte 0; datos/ultra9/correccion y datos/i5\_10400}
  \grafica[49mm]{tandas.pdf}
  \begin{idea}Todas las ejecuciones terminaron con código 0. Cada punto se midió en un proceso nuevo y en orden aleatorio.\end{idea}
  \note{Este es el volumen de trabajo. Las dos tandas más grandes son las de los binarios oficiales, con y sin Docker, con unas cuatro mil mediciones cada una.

  Los controles son pruebas pequeñas diseñadas para aislar una causa. Por ejemplo, el mismo código dentro y fuera de Docker. Las veremos cuando toque.

  Dos detalles de método: cada punto se midió en un proceso nuevo, y el orden de los puntos se sorteó. Así, si la máquina se calienta o algo la interrumpe, el efecto se reparte entre todos los algoritmos y no perjudica siempre al mismo.}
\end{frame}

\begin{frame}{El flujo de trabajo, de principio a fin}
  \fuente{replica/completa: correr\_docker.sh, preparar.sh, correr\_nativo.sh, baterias.py, analizar.py; codigo/ejecutar\_pruebas.sh}
  \vspace{2mm}
  \centering
  \begin{tikzpicture}[
      caja/.style={draw=linea,fill=suave,rounded corners=1.5pt,align=center,text width=23.5mm,minimum height=13mm,font=\scriptsize,inner sep=1.5pt},
      clave/.style={caja,draw=marino,line width=0.6pt},
      flecha/.style={-{Stealth[length=2mm]},color=marino,line width=0.6pt},
      rot/.style={font=\fontsize{6}{7}\selectfont\bfseries\color{marino},anchor=west}]
    \node[rot] at (-1.3,1.35) {BINARIOS OFICIALES DE LOS AUTORES};
    \node[clave] (z) at (0,0) {\textbf{Zenodo}\\artefacto oficial\\MD5 comprobado};
    \node[caja,right=2.6mm of z] (c) {\textbf{Compilar}\\Docker: Ubuntu 22.04, g++-12\\Sin Docker: \codp{preparar.sh}};
    \node[caja,right=2.6mm of c] (m) {\textbf{Medir}\\\codp{baterias.py}: orden al azar, un proceso por punto};
    \node[caja,right=2.6mm of m] (v) {\textbf{CSV}\\\codp{resultados/}\\\codp{david-ultra9-285}};
    \node[clave,right=2.6mm of v] (a) {\textbf{Analizar}\\\codp{analizar.py}: figuras y \codp{tablas.md}};
    \draw[flecha] (z) -- (c); \draw[flecha] (c) -- (m); \draw[flecha] (m) -- (v); \draw[flecha] (v) -- (a);
    \node[rot] at (-1.3,-1.25) {PROGRAMA DEL EQUIPO};
    \node[clave] (b) at (0,-2.55) {\textbf{\codp{bench.cpp}}\\4 algoritmos,\\incluido Mutex};
    \node[caja,right=2.6mm of b] (e) {\textbf{Pasada rápida}\\\codp{ejecutar\_pruebas.sh}\\en el i5 y en el Ultra 9};
    \node[caja,right=2.6mm of e] (n) {\textbf{Batería completa}\\\codp{correr\_nativo.sh}\\\codp{equipo} y \codp{alineado}};
    \node[caja,right=2.6mm of n] (k) {\textbf{Controles}\\Docker, alineación, Williams, RAM};
    \node[caja,right=2.6mm of k] (d) {\textbf{Documentar}\\\codp{EXPLICACION.md}\\y esta presentación};
    \draw[flecha] (b) -- (e); \draw[flecha] (e) -- (n); \draw[flecha] (n) -- (k); \draw[flecha] (k) -- (d);
    \draw[flecha] (a.south) -- (d.north);
    \draw[flecha,dashed] (n.north) -- (v.south);
  \end{tikzpicture}
  \vspace{3mm}
  \begin{idea}Todo se reanuda solo: si una medición se corta, se vuelve a lanzar la misma orden y salta lo ya medido.\end{idea}
  \note{Este es el flujo completo, en dos carriles.

  Arriba, los binarios oficiales. Se bajan de Zenodo, se comprueba que el archivo no cambió con su MD5, se compilan, se miden con un guion propio, baterias.py, que sortea el orden y lanza un proceso por punto, y todo acaba en archivos CSV. De ahí, analizar.py saca las gráficas y las tablas.

  Abajo, el programa del equipo, bench.cpp. Se corrió primero en una pasada rápida en las dos máquinas, después en una batería completa en el Ultra 9, y luego vinieron los controles.

  Todo es reanudable: si algo se corta, se relanza la misma orden y continúa donde se quedó.}
\end{frame}

\begin{frame}{Paso 1 · El artefacto oficial, descargado de Zenodo}
  \fuente{EXPLICACION.md, Partes 4.2 y Referencias; replica/completa/correr\_docker.sh}
  \begin{columns}[T,onlytextwidth]
    \column{0.55\textwidth}
    \small
    \textbf{\color{marino}Qué es}\par
    El paquete que publicaron los autores junto al paper: el código de los cinco algoritmos y los dos programas de medida.\par\vspace{2mm}
    \footnotesize
    \begin{tabular}{@{}G{17mm}p{50mm}@{}}
      DOI & \cod{10.5281/zenodo.11775182}\\
      MD5 & \codp{e7e03425dfd426385a769788e5460fcd}\\
      Programas & \cod{bench\_stack} (rendimiento) y el de latencia\\
      Rivales & Folly en el commit que fija el paper; Anderson en ParlayLib; Williams con \cod{boost::atomic}\\
      Fijado por nosotros & Google Benchmark v1.8.3 y Google Test v1.14.0
    \end{tabular}
    \column{0.41\textwidth}
    \begin{idea}\textbf{Por qué así:} se mide el mismo código que midieron los autores. Una reimplementación propia podría ganar o perder por cómo está escrita, no por el algoritmo.\end{idea}
    \begin{ojo}El artefacto pide la rama \cod{main} del día de Google Benchmark y Google Test, que ya no es la misma. Por eso se fijaron dos versiones concretas.\end{ojo}
  \end{columns}
  \note{El primer paso fue conseguir el código exacto de los autores. Lo publicaron en Zenodo, un repositorio científico, con un identificador permanente.

  Lo bajamos y comprobamos su MD5, una huella del archivo. Si coincide, sabemos que es exactamente lo que publicaron, sin un byte cambiado.

  ¿Por qué no programarlo nosotros? Porque entonces mediríamos nuestra forma de programarlo. Usando su artefacto, medimos su código, el de Folly y el de Anderson tal como ellos los usaron.

  Un detalle práctico: el artefacto pide la versión del día de dos bibliotecas de pruebas. Esa versión ya cambió, así que fijamos dos versiones concretas para que se pueda repetir.}
\end{frame}

\begin{frame}{Paso 2 · Compilar y medir dentro de Docker}
  \fuente{replica/completa/correr\_docker.sh y Dockerfile; EXPLICACION.md, Parte 9}
  \begin{columns}[T,onlytextwidth]
    \column{0.5\textwidth}
    \small
    \textbf{\color{marino}Qué hace \cod{correr\_docker.sh}}
    \begin{enumerate}\footnotesize
      \item Baja el artefacto y comprueba el MD5.
      \item Construye una imagen con Ubuntu 22.04 y g++-12, la misma base que usa el artefacto.
      \item Compila los binarios oficiales dentro.
      \item Lanza las baterías: unas 3 horas.
    \end{enumerate}
    \vspace{1mm}
    \begin{tcolorbox}[colback=tinta,colupper=white,boxrule=0pt,arc=1.5pt,left=4pt,top=2pt,bottom=2pt,fontupper=\scriptsize\ttfamily]
      cd replica/completa\\
      bash correr\_docker.sh
    \end{tcolorbox}
    \column{0.46\textwidth}
    \begin{porque}[: Docker]
      Docker no es una máquina virtual. El programa corre sobre el mismo núcleo de Linux y en la misma CPU; solo cambia lo que ve a su alrededor: bibliotecas, compilador y archivos. Sirve para tener el entorno de los autores sin tocar el sistema.
    \end{porque}
    \begin{ojo}Que Docker no altera los números había que demostrarlo. Se hizo en la parte 4.\end{ojo}
  \end{columns}
  \note{El segundo paso fue compilar y medir dentro de Docker. Un solo guion lo hace todo: baja, comprueba, construye la imagen, compila y mide. Tarda unas tres horas.

  ¿Por qué Docker? Porque nos da el mismo entorno que usaron los autores, Ubuntu 22.04 con g++-12, sin instalar nada en el sistema.

  Docker no es una máquina virtual: no hay un procesador simulado. El programa corre en la misma CPU y con el mismo núcleo de Linux. Es como correr en la misma pista con un dorsal puesto.

  Aun así, que Docker no cambie los números no se puede suponer. Lo comprobamos midiendo, y lo veremos en la parte 4.}
\end{frame}

\begin{frame}{La receta oficial, sacada del artefacto}
  \fuente{EXPLICACION.md, Parte 4.2; artefacto de Zenodo}
  \small
  \begin{columns}[T,onlytextwidth]
    \column{0.6\textwidth}
    \footnotesize
    \begin{tabular}{@{}G{20mm}p{62mm}@{}}
      \toprule
      Orden & \cod{bench\_stack -r 10 -s n -u u -i 7 --peek=true --threads=T --alg=A}\\
      Algoritmos & \cod{ivo} (Autores), \cod{mine} (Anderson), \cod{folly}, \cod{jss-free} (Williams), \cod{gnu} (Libstdc++)\\
      Asignador & \cod{jemalloc}, cargado con \cod{LD\_PRELOAD}\\
      Memoria & \cod{numactl -i all}: repartida entre los nodos\\
      Hilos & 1, 18, 36, 54, 72, 108, 144, 180, 216\\
      Tamaño de pila & 20 elementos + T/n\\
      Resultado & Mediana de 7 ejecuciones de 10 s\\
      \bottomrule
    \end{tabular}
    \column{0.36\textwidth}
    \textbf{\color{marino}Qué se mide}\par\vspace{1mm}
    \footnotesize
    \textbf{n} pilas concurrentes. En cada operación se elige una pila al azar: con probabilidad \textbf{u}\,\% se saca un elemento y se mete en otra; si no, se lee la cima.\par\vspace{2mm}
    \begin{idea}\footnotesize \textbf{n} pequeño = mucha gente en poco sitio.\\ \textbf{u} grande = mucha escritura.\end{idea}
  \end{columns}
  \note{Esta receta no aparece en el paper: la sacamos del propio artefacto, leyendo cómo lanzaban las pruebas.

  El programa crea n pilas. En cada operación elige una al azar; con probabilidad u por ciento mueve un elemento de una pila a otra, y si no, solo lee la cima.

  Para leer todo lo que viene, basta recordar dos cosas. Con n pequeño, todos los hilos se pelean por las mismas pilas: eso es contención. Con u grande, hay muchas escrituras. Con n igual a un millón, casi nunca coinciden dos hilos.

  Cada punto del paper es la mediana de siete ejecuciones de diez segundos.}
\end{frame}

\begin{frame}{Adaptaciones a una máquina de 24 núcleos, y por qué}
  \fuente{EXPLICACION.md, Parte 4.3}
  \footnotesize
  \begin{tabularx}{\linewidth}{@{}p{30mm}p{38mm}X@{}}
    \toprule
    \textbf{Paper} & \textbf{Réplica} & \textbf{Por qué}\\
    \midrule
    Hilos 1 a 216 & 1, 3, 6, 9, 12, 18, 24, 30, 36 & La lista del paper dividida entre 6: 24 es «toda la máquina», como 144; 36 es 1,5 veces, como 216\\
    7 × 10 s en un proceso & 5 rondas de 2 s, un proceso nuevo por punto, en orden aleatorio & La figura 5 pasa de unas 5 h a poco más de 1 h, y el calor o una interrupción afectan a todos por igual\\
    p = 56 en latencia & p = 1, 24 y 56 & 24 son todos los núcleos reales; con 56 hay más hilos que núcleos\\
    Williams fuera de la latencia & Williams medido, compilado con \cod{-mcx16} & El paper lo omite «porque su latencia casi siempre era mayor»: se comprobó en vez de creerlo\\
    \bottomrule
  \end{tabularx}
  \vspace{2mm}
  \begin{ojo}Recortar puede sesgar. Por eso se repitieron 30 puntos con la receta exacta de 10 s × 7: lo veremos en la parte 10.\end{ojo}
  \note{Hubo que adaptar la receta a una máquina más pequeña. Cuatro cambios, cada uno con su razón.

  Los hilos: dividimos la lista del paper entre seis. Así se conserva la proporción: 24 hilos es toda la máquina, como 144 allí.

  La duración: cinco rondas de dos segundos en vez de siete de diez. Cada punto en un proceso nuevo y en orden aleatorio. Pasamos de cinco horas a una.

  La latencia: medimos con 1, 24 y 56 hilos. Y medimos también a Williams, que el paper deja fuera. Al hacerlo encontramos un problema de compilación que veremos en la parte 7.

  Recortar tiene un precio, y lo medimos: está en la parte 10.}
\end{frame}

\begin{frame}{Fijar cada hilo a un núcleo sin tocar el binario}
  \fuente{replica/completa/herramientas/fijar\_hilos.c; EXPLICACION.md, Parte 5.4}
  \begin{columns}[T,onlytextwidth]
    \column{0.47\textwidth}
    \small
    \textbf{\color{marino}Cómo funciona}\par\vspace{1mm}
    \footnotesize
    \cod{fijar\_hilos.c} es una librería pequeña que se carga delante del programa con \cod{LD\_PRELOAD}. Intercepta la creación de cada hilo y lo ata al núcleo que le toca.\par\vspace{2mm}
    \begin{itemize}
      \item El binario oficial \textbf{no se modifica}.
      \item Se comprobó leyendo el sistema: cada hilo estaba en su núcleo.
    \end{itemize}
    \vspace{2mm}
    \begin{idea}\footnotesize El experimento limpio: \textbf{los mismos 4 hilos}, del mismo tipo de núcleo, cada vez más lejos. Solo cambia la distancia.\end{idea}
    \column{0.49\textwidth}
    \footnotesize
    \begin{tabular}{@{}p{20mm}p{44mm}@{}}
      \toprule
      \textbf{Colocación} & \textbf{Qué imita}\\
      \midrule
      \cod{E4-compacto} & 4 hilos en un clúster: todos en el mismo «socket»\\
      \cod{E4-disperso} & uno en cada clúster: repartidos por los 4 «sockets»\\
      \cod{P4} & 4 núcleos P: la máxima distancia\\
      \cod{todos24} & toda la máquina, con los hilos fijados\\
      \cod{sistema} & sin fijar: lo decide Linux, como en el paper\\
      \bottomrule
    \end{tabular}
  \end{columns}
  \note{Para simular cuatro sockets necesitamos decidir en qué núcleo exacto corre cada hilo. A eso se le llama afinidad.

  Lo hicimos sin tocar el binario oficial. Una librería pequeña se carga delante del programa, intercepta el momento en que se crea cada hilo y lo ata a un núcleo. Comprobamos después que cada hilo estaba donde debía.

  El experimento más limpio es con cuatro hilos: compacto, todos en un mismo grupo; disperso, uno en cada grupo; y en núcleos P, lo más lejos posible. Mismo número de hilos y mismo tipo de núcleo: lo único que cambia es la distancia. Es un experimento de dosis y respuesta.}
\end{frame}

\begin{frame}{Antes de simular sockets, medir la distancia real entre núcleos}
  \fuente{datos/ultra9/topologia/matriz\_nucleos.csv (herramientas/latencia\_nucleos.cpp); EXPLICACION.md, Partes 5.2 y 5.3}
  \grafica[47mm]{topologia.pdf}
  \begin{idea}Los 4 clústeres E se comportan como sockets en miniatura: cerca dentro y más lejos fuera. Pero el castigo es de 1,35×, frente a unas 3× entre sockets Xeon (dato indirecto, de máquinas parecidas).\end{idea}
  \note{Antes de suponer que los grupos de núcleos se parecen a sockets, lo medimos. Un pequeño programa lanza una pelota entre cada par de núcleos y cronometra cuánto tarda en volver.

  A la izquierda, el mapa: más oscuro es más lento. Los cuadros claros de la diagonal son los cuatro clústeres de núcleos E: dentro de cada uno, los viajes son rápidos.

  A la derecha, las medianas. Dentro del clúster, 35 nanosegundos; entre clústeres, 47. La sorpresa: los núcleos P, los más rápidos calculando, son los más lejanos entre sí, con 86 nanosegundos.

  En los Xeon de esa generación, salir del socket cuesta unas tres veces más. Aquí, 1,35. La miniatura tiene la forma correcta, pero el castigo es más suave.}
\end{frame}

\begin{frame}{Paso 3 · Sin Docker, primero con el programa del equipo}
  \fuente{EXPLICACION.md, Partes 6.6 y 8; replica/completa/correr\_nativo.sh}
  \begin{columns}[T,onlytextwidth]
    \column{0.52\textwidth}
    \small
    \textbf{\color{marino}Por qué medir sin Docker}\par\vspace{1mm}
    \footnotesize Para descartar que el contenedor cambiara los resultados. Si dentro y fuera dan lo mismo, Docker queda absuelto.\par\vspace{2mm}
    \small\textbf{\color{marino}Por qué se empezó por \cod{bench.cpp}}\par\vspace{1mm}
    \footnotesize Los binarios oficiales necesitan compilar Folly y ParlayLib fuera del contenedor. Al principio no compilaban: el sistema de permisos del asistente bloqueó \cod{preparar.sh}, y lanzado a mano se paraba al configurar \cod{bench\_stack}.
    \column{0.44\textwidth}
    \begin{tcolorbox}[colback=tinta,colupper=white,boxrule=0pt,arc=1.5pt,left=4pt,top=2pt,bottom=2pt,fontupper=\scriptsize\ttfamily]
      cd replica/completa\\
      SOLO\_DESCARGA=1 bash preparar.sh\\
      bash correr\_nativo.sh equipo
    \end{tcolorbox}
    \footnotesize
    \begin{itemize}
      \item Autores, Libstdc++, Williams y Mutex; \textbf{Folly y Anderson no están}.
      \item Se le añadió jemalloc con \cod{LD\_PRELOAD} y n = 10⁶, sin tocar su código.
      \item 2\,976 mediciones, unas 1,5 horas.
    \end{itemize}
  \end{columns}
  \vspace{2mm}
  \begin{ojo}Aquí cambian dos cosas a la vez, el entorno y el programa. Separarlas exigió los controles de las partes 4 y 6.\end{ojo}
  \note{El tercer paso fue medir sin Docker. La idea era sencilla: si dentro y fuera da lo mismo, el contenedor no influye.

  Empezamos con el programa del equipo, bench.cpp, porque los binarios oficiales no compilaban fuera del contenedor. Necesitan compilar Folly y ParlayLib, y el sistema de permisos del asistente bloqueó ese paso por tratarse de código descargado. Lanzado a mano, se paraba a mitad.

  bench.cpp tiene cuatro algoritmos: le faltan Folly y Anderson. Le añadimos jemalloc y el escenario de un millón de pilas, sin tocar su código.

  Ojo: aquí cambiaban a la vez el entorno y el programa. Si salía distinto, no sabíamos a quién culpar. Por eso hicieron falta controles.}
\end{frame}

\begin{frame}{Paso 4 · Sin Docker con los binarios oficiales: la tanda nocturna}
  \fuente{EXPLICACION.md, Partes 8 y 9; replica/completa/pendientes.sh}
  \begin{columns}[T,onlytextwidth]
    \column{0.5\textwidth}
    \small
    \textbf{\color{marino}El fallo, y su arreglo}\par\vspace{1mm}
    \footnotesize El archivo de CMake de \cod{fmt} sube 4 carpetas para hallar su raíz, pero el enlace de \cod{preparar.sh} estaba a 3 carpetas de \cod{install/fmt}: buscaba \cod{install/lib/...} y no existía. Se arregló con dos enlaces.\par\vspace{2.5mm}
    \small\textbf{\color{marino}Qué hizo \cod{pendientes.sh}}
    \begin{enumerate}\footnotesize
      \item Compilar los binarios oficiales fuera de Docker.
      \item Medir los 5 algoritmos: figura 5, afinidad y pares (10 rondas).
      \item Latencia con 15 rondas en vez de 5.
      \item Repetir \cod{bench.cpp} con las pilas alineadas.
      \item Regenerar tablas y gráficas.
    \end{enumerate}
    \column{0.46\textwidth}
    \vspace{1mm}
    \cifra{4 h 44 min}{noche del 6 al 7 de noviembre de 2026,\\de 23:07 a 03:51}\par\vspace{4mm}
    \cifra{0}{fallos; \cod{comprobar\_final.py} confirmó\\que no faltaba ninguna fila}
  \end{columns}
  \note{El cuarto paso fue repetir con los binarios oficiales, ya fuera de Docker. Primero hubo que entender por qué no compilaban.

  El culpable era fmt, una biblioteca de formato que usa Folly. Su archivo de configuración sube cuatro carpetas para encontrar su raíz, y nuestro enlace estaba a tres. Buscaba una carpeta que no existía. Dos enlaces más y compiló.

  Después, un solo guion, pendientes.sh, hizo todo de un tirón durante la noche: compilar, medir los cinco algoritmos, repetir la latencia con quince rondas y medir la versión alineada del programa del equipo. Tardó 4 horas y 44 minutos, sin un solo fallo. Otro guion comprobó al final que no faltaba nada.}
\end{frame}

\begin{frame}{Dentro y fuera de Docker cambia más que Docker}
  \fuente{EXPLICACION.md, Parte 6.6}
  \begin{columns}[T,onlytextwidth]
    \column{0.5\textwidth}
    \footnotesize
    \begin{tabular}{@{}lll@{}}
      \toprule
      & \textbf{Con Docker} & \textbf{Sin Docker}\\
      \midrule
      g++-12 & 12.3.0 & 12.4.0\\
      glibc & 2.35 & 2.39\\
      Boost & 1.74 & 1.83\\
      jemalloc & 5.2.1 & 5.3.0\\
      Folly & \cod{getdeps.py} & CMake directo\\
      \bottomrule
    \end{tabular}
    \column{0.46\textwidth}
    \begin{idea}\footnotesize Si un resultado cambia al salir de Docker, no se puede culpar a Docker sin más: cambian cinco piezas a la vez.\end{idea}
    \begin{idea}\footnotesize Y si nada cambia, es que esas cinco piezas juntas tampoco importan aquí.\end{idea}
  \end{columns}
  \vspace{3mm}
  \begin{porque}[: por qué importa Boost]
    El Williams del artefacto usa \cod{boost::atomic}. Un cambio de versión de Boost es el primer sospechoso de la única diferencia grande entre las dos tandas, Williams entre clústeres (parte 9). Es una sospecha: no se midió.
  \end{porque}
  \note{Un detalle que es fácil pasar por alto. Salir de Docker no cambia solo Docker. Cambian el compilador, la biblioteca de C, Boost, jemalloc y la forma de compilar Folly.

  Esto tiene dos lecturas. Si algo cambia, no podemos culpar a Docker sin más, porque hay cinco sospechosos. Y si nada cambia, sabemos que esas cinco piezas tampoco importan.

  Lo de Boost lo retomaremos: el Williams del artefacto usa Boost, y es el único algoritmo con una diferencia grande entre las dos tandas. Pero es una sospecha que no medimos.}
\end{frame}

\begin{frame}{El programa del equipo no es el programa oficial}
  \fuente{EXPLICACION.md, Parte 6.6}
  \footnotesize
  \begin{tabularx}{\linewidth}{@{}G{30mm}LL@{}}
    \toprule
    & \textbf{\cod{bench\_stack}, oficial} & \textbf{\cod{codigo/bench.cpp}, del equipo}\\
    \midrule
    Algoritmos & Los 5 del paper & Autores, Libstdc++, Williams y Mutex\\
    Pilas en memoria & Cada una en su sitio: \cod{alignas(128)} & Seguidas, sin alinear\\
    Nodos por pila & 20 + T/n & 16, o 2 si n > 10\\
    Generador de azar & Uno rápido & \cod{mt19937\_64}, más lento\\
    Compilación & \cod{-O3} & \cod{-O2}\\
    Williams & Con \cod{boost::atomic} & Con \cod{std::atomic}\\
    Autores & Versión del artefacto: la cuenta sube & Versión publicada: la cuenta baja\\
    \bottomrule
  \end{tabularx}
  \vspace{2mm}
  \begin{idea}Son dos programas escritos por personas distintas. Cuando coinciden, el resultado es sólido. Cuando no, la diferencia está en el programa, y hay que encontrarla.\end{idea}
  \note{bench.cpp lo escribió el equipo como comprobación independiente. No es el programa oficial, y conviene saber en qué se diferencian.

  La diferencia que más importará es la segunda fila: el oficial coloca cada pila en su propio hueco de memoria, y el nuestro las pone seguidas. Parece un detalle y nos costó un factor tres. Lo veremos en la parte 6.

  También cambian el número de nodos, el generador de azar, el nivel de optimización y la versión de Williams y de Autores.

  La idea de fondo: dos programas distintos que cuentan la misma historia dan mucha confianza. Y donde no coinciden, hay algo que aprender.}
\end{frame}

\begin{frame}{El i5-10400: cómo se corrieron sus pruebas}
  \fuente{codigo/ejecutar\_pruebas.sh y correr\_tsan.sh; datos/i5\_10400/00\_entorno.txt}
  \begin{columns}[T,onlytextwidth]
    \column{0.5\textwidth}
    \small
    \textbf{\color{marino}Entorno}\par\vspace{1mm}
    \footnotesize
    \begin{tabular}{@{}G{17mm}p{47mm}@{}}
      Sistema & WSL2 sobre Windows (Linux 6.18), Ubuntu 24.04\\
      Compilador & g++ 13.3\\
      Hilos & 12 (6 núcleos con SMT)\\
    \end{tabular}\par\vspace{2mm}
    \small\textbf{\color{marino}Qué hace \cod{ejecutar\_pruebas.sh}}
    \begin{enumerate}\footnotesize
      \item Demo del código de los autores.
      \item ASan con 8, 12 y 16 hilos; TSan con \cod{setarch -R}.
      \item El boceto y dos variantes rotas.
      \item \cod{bench.cpp}: rendimiento, 1 s por punto, hilos 1 a 18.
      \item \cod{bench.cpp}: latencia, 100\,000 muestras, p = 1 y 12.
    \end{enumerate}
    \column{0.46\textwidth}
    \begin{tcolorbox}[colback=tinta,colupper=white,boxrule=0pt,arc=1.5pt,left=4pt,top=2pt,bottom=2pt,fontupper=\scriptsize\ttfamily]
      cd codigo\\
      bash ejecutar\_pruebas.sh \textbackslash\\
      \ \ ../replica/resultados/<nombre>-<cpu>
    \end{tcolorbox}
    \begin{idea}\footnotesize El mismo guion se corrió en el Ultra 9 (\cod{correccion/}): mismos parámetros, misma versión del compilador.\end{idea}
    \begin{ojo}Una sola pasada de 1 s por punto: solo valen las diferencias grandes.\end{ojo}
  \end{columns}
  \note{Ahora el i5-10400. Corre Linux dentro de Windows, con WSL2, y el compilador g++ 13.3. Tiene 6 núcleos y cada uno ejecuta dos hilos a la vez, así que el sistema ve 12.

  Todo se lanzó con un solo guion del equipo, ejecutar\_pruebas.sh. Hace cinco cosas: la demo del código de los autores, los detectores de errores ASan y TSan, el boceto con sus variantes rotas, y una pasada rápida de rendimiento y de latencia con bench.cpp.

  TSan necesitó un truco: desactivar la aleatorización de direcciones con setarch, porque en estos núcleos de Linux fallaba al arrancar.

  Lo importante: el mismo guion se corrió en el Ultra 9. Por eso podemos compararlas. Pero es una sola pasada de un segundo por punto: solo nos fiaremos de las diferencias grandes.}
\end{frame}

\begin{frame}{Controles: un experimento pequeño para cada duda}
  \fuente{replica/completa: controles.sh, control\_williams.sh, herramientas/; EXPLICACION.md, Partes 6.1 a 6.7}
  \footnotesize
  \begin{tabularx}{\linewidth}{@{}p{38mm}Xp{17mm}@{}}
    \toprule
    \textbf{Control} & \textbf{Qué aísla} & \textbf{Dónde se ve}\\
    \midrule
    Mismo código dentro y fuera de Docker & Docker por sí solo: 48 puntos alternados, 3 rondas & Parte 4\\
    \cod{alignas(128)} en una copia de \cod{bench.cpp} & La colocación de las pilas en memoria & Parte 6\\
    Williams con \cod{std::atomic} y con \cod{boost::atomic} & La biblioteca atómica de Williams & Parte 9\\
    Contadores de la controladora de memoria & El tráfico real de RAM: \cod{medir\_memoria.py}, \cod{alinear\_memoria.py} & Partes 3 y 6\\
    Detectores en los rivales & Que ningún rival gane usando memoria liberada & Parte 2\\
    Contador en una copia del código de los autores & Los intentos reales de la recarga & Parte 2\\
    Receta exacta, 10 s × 7 & El efecto de recortar la duración & Parte 10\\
    Valores del paper sacados del PDF vectorial & Comparar número contra número (\cod{extraer\_fig5.py}) & Partes 3 y 7\\
    \bottomrule
  \end{tabularx}
  \note{Cada vez que un resultado no cuadraba, diseñamos un control: un experimento pequeño que cambia una sola cosa.

  Por ejemplo, para saber si Docker influye, compilamos el mismo código dentro y fuera y lo medimos alternando. Para saber si la colocación de las pilas en memoria importaba, añadimos una sola palabra en una copia del programa.

  Una mención especial a la última fila. Los valores del paper no los leímos a ojo de la gráfica. Las figuras del PDF son vectoriales, así que un guion saca la posición exacta de cada marcador y la convierte a Mop/s. Así comparamos número contra número.}
\end{frame}

\begin{frame}{Calidad de las mediciones y tres reglas para leerlas}
  \fuente{EXPLICACION.md, Partes 6 y 8}
  \begin{columns}[T,onlytextwidth]
    \column{0.47\textwidth}
    \small\textbf{\color{marino}Calidad}\par\vspace{1mm}
    \footnotesize
    \begin{itemize}
      \item 0 ejecuciones fallidas en todas las tandas.
      \item Con Docker: temperatura mediana de 75\,°C; solo 5 de 3\,720 llegaron a 90\,°C o más.
      \item Sin Docker: 26 de 9\,796 a 90\,°C o más; máximo, 99\,°C. Repartidas, no se descartaron.
      \item Dispersión típica entre rondas: 1\,\%; en los puntos más ruidosos, 8\,\%.
      \item PC de escritorio sin aislar y con la frecuencia sin fijar: el orden aleatorio reparte las interferencias.
    \end{itemize}
    \column{0.49\textwidth}
    \small\textbf{\color{marino}Tres reglas de lectura}\par\vspace{1mm}
    \footnotesize
    \begin{enumerate}
      \item \textbf{Se comparan algoritmos, no máquinas.} La CPU es casi diez años más moderna que los Xeon del paper.
      \item \textbf{«0,50×» es «rinde la mitad que Autores».} Así se compara sin que importe la máquina.
      \item \textbf{Una diferencia pequeña es un empate.} Recortar la duración movió algunos puntos hasta un 21\,\%: solo cuentan las diferencias de 2 o 3 veces.
    \end{enumerate}
  \end{columns}
  \note{Antes de entrar en los resultados, la calidad de los datos y tres reglas para leerlos.

  Calidad: ninguna ejecución falló. El procesador se calentó, como es normal, pero muy pocas mediciones pasaron de 90 grados, y estaban repartidas entre escenarios, así que no se descartaron. La variación típica entre rondas fue del uno por ciento.

  Las reglas. Primera: comparamos algoritmos, no máquinas. Segunda: casi todo lo expresamos respecto a Autores, el algoritmo del paper; 0,50 significa que rinde la mitad. Tercera, la más importante: una diferencia pequeña no es una victoria. Lo justificaremos en la parte 10.}
\end{frame}
